% Part: model-theory
% Chapter: interpolation
% Section: definability

\documentclass[../../../include/open-logic-section]{subfiles}

\begin{document}

\olfileid{mod}{int}{def}

\olsection{مبرهنة قابلية التعريف}

تثمر مبرهنة الاستيفاء نتائج مهمة، منها مبرهنة بث في قابلية التعريف.
فالعلاقة~${\OLId{latin-looped}{R}}$ ذات ${\OLId{latin-ordinary}{n}}$ مواضع يصح تعريفها بإيراد !!a{formula}~$!C$ ذات
${\OLId{latin-ordinary}{n}}$~!!{variable}s حرة، خالية من~${\OLId{latin-looped}{R}}$؛ فيكون ذلك تعريفًا \emph{صريحًا} لـ${\OLId{latin-looped}{R}}$
بدلالة~$!C$. وعلى هذا يقال إن نظرية~$\Sigma({\OLId{latin-looped}{P}})$ في !!a{language}
تشتمل على الـ!!{predicate}~${\OLId{latin-looped}{P}}$ ذي ${\OLId{latin-ordinary}{n}}$ مواضع تعرّف~${\OLId{latin-looped}{P}}$ صراحة إذا
اشتملت على تعريف صريح محرر في تلك اللغة، أو استلزمته وإن لم يرد فيها، أي:
\[
\Sigma({\OLId{latin-looped}{P}}) \Entails \lforall[{\OLId{latin-ordinary}{x}}_{\OLMathNumeral{1}}][\dots
  \lforall[{\OLId{latin-ordinary}{x}}_{\OLId{latin-ordinary}{n}}][(\Atom{P}{x_1,\dots, x_n} \liff !C({\OLId{latin-ordinary}{x}}_{\OLMathNumeral{1}}, \dots,
    {\OLId{latin-ordinary}{x}}_{\OLId{latin-ordinary}{n}}))]].
\]
وليس التصريح وحده سبيلًا إلى تعريف العلاقة، إذا أريد بالتعريف تعيينها
تعيينًا تامًا. فقد تقتضي النظرية أن يتعين تفسير~${\OLId{latin-looped}{P}}$ في كل نموذج بمجرد تعيين بقية الـ!!{language}.
ومؤدى مبرهنة قابلية التعريف أن النظرية متى عيّنت تفسير~${\OLId{latin-looped}{P}}$
على هذا الوجه، أي \emph{عرّفت ${\OLId{latin-looped}{P}}$ ضمنيًا}، لزم أن تعرّفه صراحة أيضًا.

\begin{defn}
لتكن $\Lang{L}$~!!a{language} خالية من الـ!!{predicate}~${\OLId{latin-looped}{P}}$.
تعرّف مجموعة~$\Sigma({\OLId{latin-looped}{P}})$ من !!{sentence}s~$\Lang{L} \cup \{{\OLId{latin-looped}{P}}\}$
المحمول~${\OLId{latin-looped}{P}}$ \emph{تعريفًا صريحًا} إذا وفقط إذا وجدت
!!a{formula}~$!C({\OLId{latin-ordinary}{x}}_{\OLMathNumeral{1}},\dots,{\OLId{latin-ordinary}{x}}_{\OLId{latin-ordinary}{n}})$ من $\Lang{L}$ بحيث
\[
\Sigma({\OLId{latin-looped}{P}}) \Entails \lforall[{\OLId{latin-ordinary}{x}}_{\OLMathNumeral{1}}][\dots
  \lforall[{\OLId{latin-ordinary}{x}}_{\OLId{latin-ordinary}{n}}][(\Atom{P}{x_1,\dots, x_n} \liff !C({\OLId{latin-ordinary}{x}}_{\OLMathNumeral{1}}, \dots,
    {\OLId{latin-ordinary}{x}}_{\OLId{latin-ordinary}{n}}))]].
\]
\end{defn}

\begin{defn}
لتكن $\Lang{L}$~!!a{language} لا يرد فيها الـ!!{predicate}s~${\OLId{latin-looped}{P}}$ و
${\OLId{latin-looped}{P}}'$. تعرّف مجموعة~$\Sigma({\OLId{latin-looped}{P}})$ من !!{sentence}s~$\Lang{L} \cup \{{\OLId{latin-looped}{P}}\}$
المحمول~${\OLId{latin-looped}{P}}$ \emph{تعريفًا ضمنيًا} إذا وفقط إذا كان
\[
\Sigma({\OLId{latin-looped}{P}}) \cup \Sigma({\OLId{latin-looped}{P}}') \Entails \lforall[{\OLId{latin-ordinary}{x}}_{\OLMathNumeral{1}}][\dots
  \lforall[{\OLId{latin-ordinary}{x}}_{\OLId{latin-ordinary}{n}}][(\Atom{P}{x_1,\dots, x_n} \liff \Atom{P'}{x_1,\dots,
      x_n})]],
\]
حيث $\Sigma({\OLId{latin-looped}{P}}')$ ناتجة من إحلال ${\OLId{latin-looped}{P}}'$ محل ${\OLId{latin-looped}{P}}$ بانتظام في $\Sigma({\OLId{latin-looped}{P}})$.
\end{defn}

وتفسير هذا الشرط: أنه مهما أخذنا نموذجًا~$\Struct{M}$ وعلاقتين
${\OLId{latin-looped}{R}},{\OLId{latin-looped}{R}}' \subseteq \Domain{{\OLId{latin-looped}{M}}}^{\OLId{latin-ordinary}{n}}$، إذا كان كل من
$\Expan{{\OLId{latin-looped}{M}}}{{\OLId{latin-looped}{R}}} \Entails \Sigma({\OLId{latin-looped}{P}})$ و
$\Expan{{\OLId{latin-looped}{M}}}{{\OLId{latin-looped}{R}}'} \Entails \Sigma({\OLId{latin-looped}{P}}')$، فإن ${\OLId{latin-looped}{R}}={\OLId{latin-looped}{R}}'$؛ حيث
$\Expan{{\OLId{latin-looped}{M}}}{{\OLId{latin-looped}{R}}}$ هي الـ!!{structure}~$\Struct{M'}$ لتوسيع $\Lang{L}$
إلى $\Lang{L} \cup \{{\OLId{latin-looped}{P}}\}$ بحيث $\Assign{{\OLId{latin-looped}{P}}}{{\OLId{latin-looped}{M}}'}={\OLId{latin-looped}{R}}$، وبالمثل لـ
$\Expan{{\OLId{latin-looped}{M}}}{{\OLId{latin-looped}{R}}'}$.

\begin{thm}[مبرهنة بث في قابلية التعريف]
تعرّف مجموعة~$\Sigma({\OLId{latin-looped}{P}})$ من !!{sentence}s~$\Lang{L} \cup \{{\OLId{latin-looped}{P}}\}$
المحمول~${\OLId{latin-looped}{P}}$ تعريفًا ضمنيًا إذا وفقط إذا كانت $\Sigma({\OLId{latin-looped}{P}})$ تعرّف~${\OLId{latin-looped}{P}}$
تعريفًا صريحًا.
\end{thm}

\begin{proof}
إذا كانت $\Sigma({\OLId{latin-looped}{P}})$ تعرّف ${\OLId{latin-looped}{P}}$ تعريفًا صريحًا، كان
\begin{align*}
  \Sigma({\OLId{latin-looped}{P}}) & \Entails & \lforall[{\OLId{latin-ordinary}{x}}_{\OLMathNumeral{1}}][\dots \lforall[{\OLId{latin-ordinary}{x}}_{\OLId{latin-ordinary}{n}}]
    [(\Atom{P}{x_1,\dots, x_n} \liff !C({\OLId{latin-ordinary}{x}}_{\OLMathNumeral{1}},\dots,{\OLId{latin-ordinary}{x}}_{\OLId{latin-ordinary}{n}}))]]\\
  \Sigma({\OLId{latin-looped}{P}}') & \Entails & \lforall[{\OLId{latin-ordinary}{x}}_{\OLMathNumeral{1}}][\dots \lforall[{\OLId{latin-ordinary}{x}}_{\OLId{latin-ordinary}{n}}]
    [(\Atom{P'}{x_1,\dots, x_n} \liff !C({\OLId{latin-ordinary}{x}}_{\OLMathNumeral{1}},\dots,{\OLId{latin-ordinary}{x}}_{\OLId{latin-ordinary}{n}}))]]
\end{align*}
فيثبت المطلوب من هذين معًا. وأما العكس، فنفرض أن $\Sigma({\OLId{latin-looped}{P}})$ تعرّف~${\OLId{latin-looped}{P}}$ تعريفًا
ضمنيًا، ونبدأ بإلحاق الـ!!{constant}s~${\OLId{latin-ordinary}{c}}_{\OLMathNumeral{1}}$، \dots،~${\OLId{latin-ordinary}{c}}_{\OLId{latin-ordinary}{n}}$ بـ
$\Lang{L}$. فيلزم حينئذ
\[
\Sigma({\OLId{latin-looped}{P}}) \cup \Sigma({\OLId{latin-looped}{P}}') \Entails
\Atom{P}{c_1, \dots, c_n} \to \Atom{P'}{c_1, \dots, c_n}.
\]
ومن التراص نحصل على مجموعتين منتهيتين
${\OLId{greek-variable}{Delta}}_{\OLMathNumeral{0}} \subseteq \Sigma({\OLId{latin-looped}{P}})$ و${\OLId{greek-variable}{Delta}}_{\OLMathNumeral{1}} \subseteq \Sigma({\OLId{latin-looped}{P}}')$ بحيث
\[
{\OLId{greek-variable}{Delta}}_{\OLMathNumeral{0}} \cup {\OLId{greek-variable}{Delta}}_{\OLMathNumeral{1}} \Entails
\Atom{P}{c_1, \dots, c_n} \to \Atom{P'}{c_1, \dots, c_n}.
\]
لتكن $!D({\OLId{latin-looped}{P}})$ اقتران جميع الـ!!{sentence}s~$!A({\OLId{latin-looped}{P}})$ بحيث إما أن
$!A({\OLId{latin-looped}{P}}) \in {\OLId{greek-variable}{Delta}}_{\OLMathNumeral{0}}$ أو أن $!A({\OLId{latin-looped}{P}}') \in {\OLId{greek-variable}{Delta}}_{\OLMathNumeral{1}}$، ولتكن $!D({\OLId{latin-looped}{P}}')$
اقتران جميع الـ!!{sentence}s~$!A({\OLId{latin-looped}{P}}')$ بحيث إما أن
$!A({\OLId{latin-looped}{P}}) \in {\OLId{greek-variable}{Delta}}_{\OLMathNumeral{0}}$ أو أن $!A({\OLId{latin-looped}{P}}') \in {\OLId{greek-variable}{Delta}}_{\OLMathNumeral{1}}$. عندئذ
$!D({\OLId{latin-looped}{P}}) \land !D({\OLId{latin-looped}{P}}') \Entails \Atom{P}{c_1,\dots,c_n} \to
\Atom{P'}{c_1,\dots,c_n}$. وإذا نقلنا أجزاء هذا الاستلزام إلى مواضعها المناسبة، اقتصر ورود كل
!!{predicate} على أحد جانبي $\Entails$:
\[
!D({\OLId{latin-looped}{P}}) \land \Atom{P}{c_1, \dots, c_n} \Entails
!D({\OLId{latin-looped}{P}}') \to \Atom{P'}{c_1, \dots, c_n}.
\]
فتقتضي مبرهنة كريغ في الاستيفاء وجود !!a{sentence}~$!C({\OLId{latin-ordinary}{c}}_{\OLMathNumeral{1}},\dots,{\OLId{latin-ordinary}{c}}_{\OLId{latin-ordinary}{n}})$ لا
تحتوي ${\OLId{latin-looped}{P}}$ أو ${\OLId{latin-looped}{P}}'$ بحيث:
\begin{align*}
  !D({\OLId{latin-looped}{P}}) \land \Atom{P}{c_1, \dots, c_n} & \Entails !C({\OLId{latin-ordinary}{c}}_{\OLMathNumeral{1}},\dots, {\OLId{latin-ordinary}{c}}_{\OLId{latin-ordinary}{n}}); \\
  !C({\OLId{latin-ordinary}{c}}_{\OLMathNumeral{1}},\dots, {\OLId{latin-ordinary}{c}}_{\OLId{latin-ordinary}{n}}) & \Entails !D({\OLId{latin-looped}{P}}') \to \Atom{P'}{c_1, \dots, c_n}.
\end{align*}
أما الاستلزام الأول فيعطي
$!D({\OLId{latin-looped}{P}}) \Entails \Atom{P}{c_1,\dots,c_n} \lif !C({\OLId{latin-ordinary}{c}}_{\OLMathNumeral{1}},\dots,{\OLId{latin-ordinary}{c}}_{\OLId{latin-ordinary}{n}})$.
وأما الثاني فنستعمل فيه أن نموذج $\Lang{L} \cup \{{\OLId{latin-looped}{P}}\}$
$\Expan{{\OLId{latin-looped}{M}}}{{\OLId{latin-looped}{R}}} \models !A({\OLId{latin-looped}{P}})$ إذا وفقط إذا كان نموذج
$\Lang{L} \cup \{{\OLId{latin-looped}{P}}'\}$ المقابل $\Expan{{\OLId{latin-looped}{M}}}{{\OLId{latin-looped}{R}}} \models !A({\OLId{latin-looped}{P}}')$؛ فينتج $!C({\OLId{latin-ordinary}{c}}_{\OLMathNumeral{1}},\dots,{\OLId{latin-ordinary}{c}}_{\OLId{latin-ordinary}{n}}) \Entails !D({\OLId{latin-looped}{P}}) \lif
\Atom{P}{c_1,\dots,c_n}$، ومنه:
\[
!D({\OLId{latin-looped}{P}}) \Entails !C({\OLId{latin-ordinary}{c}}_{\OLMathNumeral{1}},\dots,{\OLId{latin-ordinary}{c}}_{\OLId{latin-ordinary}{n}}) \to \Atom{P}{c_1,\dots, c_n}.
\]
فإذا جمعنا النتيجتين حصل
$!D({\OLId{latin-looped}{P}}) \Entails \Atom{P}{c_1,\dots,c_n} \liff
!C({\OLId{latin-ordinary}{c}}_{\OLMathNumeral{1}},\dots,{\OLId{latin-ordinary}{c}}_{\OLId{latin-ordinary}{n}})$، ثم تقتضي الرتابة والتعميم:
\[
\Sigma({\OLId{latin-looped}{P}}) \Entails
\lforall[{\OLId{latin-ordinary}{x}}_{\OLMathNumeral{1}}][\dots\lforall[{\OLId{latin-ordinary}{x}}_{\OLId{latin-ordinary}{n}}][(\Atom{P}{x_1,\dots, x_n} \liff
    !C({\OLId{latin-ordinary}{x}}_{\OLMathNumeral{1}},\dots, {\OLId{latin-ordinary}{x}}_{\OLId{latin-ordinary}{n}}))]].
\]
\end{proof}

\end{document}
